\documentclass[../../main.tex]{subfiles}% 注意这里的文件路径不能用 ./main.tex ,否则用latexmk编译子文件会报错
\graphicspath{{\subfix{./image/}}{\subfix{../../image/}}} % 指定图片引用路径,后续可以直接使用图片文件名
% 如果图片存放在上级目录,则使用 ../../image/ 作为备用路径

% 例如:
% \begin{figure}[H]
% \centering
% \includegraphics[width=0.3\textwidth]{图.png}
% \caption{}
% \label{figure:图}
% \end{figure}
% 注意:上述\label{}一定要放在\caption{}之后,否则引用图片序号会只会显示??.

\begin{document}

\section{可测集与测度}

\begin{definition}[可测集]
设 \(E \subset \mathbb{R}^n\). 若对任意的点集 \(T \subset \mathbb{R}^n\), 有
\begin{align*}
m^*(T)=m^*(T \cap E)+m^*(T \cap E^c), 
\end{align*}
则称 \(E\) 为 \textbf{Lebesgue 可测集}(或 \textbf{\(m^*\)-可测集})或$E$\textbf{可测}, 简称为\textbf{可测集}, 其中 \(T\) 称为\textbf{试验集}(这一定义可测集的等式也称为 \textbf{Carathéodory条件}). 可测集的全体称为\textbf{可测集类}, 简记为 \(\mathscr{M}\). 
\end{definition}

\begin{theorem}[集合可测的充要条件]\label{theorem:可测的充要条件}
设\(E\subset \mathbb{R}^n\), 则$E\in \mathscr{M}$的充要条件是对任一点集 \(T \subset \mathbb{R}^n\)且\(m^*(T)< + \infty\),都有
\begin{align}
m^*(T) \geqslant  m^*(T \cap E)+m^*(T \cap E^c) \label{eq:2.3}
\end{align}
成立. 
\end{theorem}
\begin{remark}
往后经常利用这个定理的充分性来证明一个集合可测.但这个定理的必要性要弱于可测集的定义.
\end{remark}
\begin{proof}
必要性由可测集的定义立得.下证充分性.
由外测度的次可加性可得
\[
m^*(T)=m^*\left( T\cap \mathbb{R} ^n \right) =m^*\left( T\cap \left( E\cup E^c \right) \right) =m^*\left( \left( T\cap E \right) \cup \left( T\cap E^c \right) \right)  \leqslant  m^*(T \cap E)+m^*(T \cap E^c)
\]
总是成立的. 又因为在 \(m^*(T)=\infty\) 时 \eqref{eq:2.3} 式总成立,故对任意的点集 \(T \subset \mathbb{R}^n\),都有
\begin{align*}
m^*(T)=m^*(T \cap E)+m^*(T \cap E^c), 
\end{align*}
即$E\in \mathscr{M}$.
\end{proof}

\begin{definition}[零测集]
外测度为零的点集称为\textbf{零测集}. 
\end{definition}
\begin{remark}
显然, \(\mathbb{R}^n\) 中由单个点组成的点集是零测集. 从而根据外测度的次可加性知道 \(\mathbb{R}^n\) 中的有理点集 \(\mathbb{Q}^n\) 是零测集.
\end{remark}

\begin{proposition}\label{proposition:零测集的基本性质}
\begin{enumerate}
\item 零测集的任一子集是零测集.

\item 零测集一定可测,即若 \(m^*(E)=0\), 则 \(E \in \mathscr{M}\). 
\end{enumerate}
\end{proposition}
\begin{proof}
\begin{enumerate}
\item 由外测度的单调性立得.

\item 事实上, 此时我们有
\begin{align*}
m^*(T \cap E)+m^*(T \cap E^c) \leqslant  m^*(E)+m^*(T)=m^*(T).
\end{align*}
再由\refthe{theorem:可测的充要条件}立得.
\end{enumerate}
\end{proof}

\begin{proposition}\label{proposition:外测度的可加性的条件}
若 \(E_1 \subset S\), \(E_2 \subset S^c\), \(S \in \mathscr{M}\), 则有
\begin{align*}
m^*(E_1 \cup E_2) = m^*(E_1) + m^*(E_2).
\end{align*}
\end{proposition}
\begin{remark}
这个命题表明:\textbf{当两个集合由一个可测集分离开时, 其外测度就有可加性}.
\end{remark}
\begin{proof}
事实上, 此时取试验集 \(T = E_1 \cup E_2\), 从 \(S\) 是可测集的定义得
\begin{align*}
m^*(E_1 \cup E_2) = m^*((E_1 \cup E_2) \cap S) + m^*((E_1 \cup E_2) \cap S^c) = m^*(E_1) + m^*(E_2).
\end{align*}
\end{proof}

\begin{corollary}\label{corollary:2个无交可测集与任意集合的并的外测度的性质}
当 \(E_1\) 与 \(E_2\) 是互不相交的可测集时, 对任一集合 \(T\) 有
\begin{align*}
m^*(T \cap (E_1 \cup E_2)) =m^*((T\cap E_1)\cup (T\cap E_2))= m^*(T \cap E_1) + m^*(T \cap E_2).
\end{align*} 
\end{corollary}
\begin{proof}
注意到$T\cap E_1\in E_1$,$T\cap E_2\in E_1^c$,而$E_1\in \mathscr{M}$,故由集合运算的性质和\refpro{proposition:外测度的可加性的条件}可知
\begin{align*}
m^*(T \cap (E_1 \cup E_2))=m^*((T\cap E_1)\cup (T\cap E_2)) = m^*(T \cap E_1) + m^*(T \cap E_2).
\end{align*} 
\end{proof}

\begin{corollary}\label{corollary:有限个无交可测集与任意集合的并的外测度的性质}
当$E_1,E_2,\cdots,E_n$是互不相交的可测集时，对任一集合$T$有
\begin{align*}
m^*\left(T\cap\bigcup_{k=1}^nE_k\right) = m^*\left(\bigcup_{k=1}^n\left(T\cap E_k\right)\right) = \sum_{k=1}^n m^*(T\cap E_k).
\end{align*}
特别地,若还有$F_k\subseteq E_k,k=1,2,\cdots,n$.则
\begin{align*}
m^*\left( \bigcup_{k=1}^n{F_k} \right) =\sum_{k=1}^n{m^*\left( F_k \right)}.
\end{align*}
\end{corollary}
\begin{proof}
当$n=1$时, 结论显然成立.假设当$n=m$时结论成立,考虑$n=m+1$的情况.由于$E_1,E_2,\cdots,E_{m+1}$皆互不相交,因此$\bigcup_{k=1}^m{E_k}$和$E_{m+1}$也互不相交.于是由集合运算的性质和\refcor{corollary:2个无交可测集与任意集合的并的外测度的性质}以及归纳假设可得
\begin{align*}
&m^*\left( T\cap \bigcup_{k=1}^{m+1}{E_k} \right) =m^*\left( \bigcup_{k=1}^{m+1}{\left( T\cap E_k \right)} \right) =m^*\left( T\cap \left( \bigcup_{k=1}^m{E_k}\cup E_{m+1} \right) \right) =m^*\left( \left( T\cap \bigcup_{k=1}^m{E_k} \right) \cup \left( T\cap E_{m+1} \right) \right) 
\\
&=m^*\left( T\cap \bigcup_{k=1}^m{E_k} \right) +m^*\left( T\cap E_{m+1} \right) =\sum_{k=1}^m{m^*\left( T\cap E_k \right)}+m^*\left( T\cap E_{m+1} \right) =\sum_{k=1}^{m+1}{m^*\left( T\cap E_k \right)}.
\end{align*}
故由数学归纳法可知结论成立.

特别地,当$F_k\subseteq E_k,k=1,2,\cdots,n$时,取$T=\bigcup_{k=1}^n{F_k}$,则
\begin{align*}
m^*\left( \bigcup_{k=1}^n{F_k} \right) &=m^*\left( \bigcup_{k=1}^n{\left( F_k\cap E_k \right)} \right) =m^*\left[ \bigcup_{k=1}^n{\left( \left( \bigcup_{i=1}^n{F_i} \right) \cap E_k \right)} \right] =\sum_{k=1}^n{m^*\left( \left( \bigcup_{i=1}^n{F_i} \right) \cap E_k \right)}
\\
&=\sum_{k=1}^n{m^*\left( F_k\cap E_k \right)}=\sum_{k=1}^n{m^*\left( F_k \right)}.
\end{align*}
\end{proof}

\begin{proposition}
设 $\{A_n\}$ 是互不相交的可测集列, $B_n\subseteq A_n$ ($n=1,2,\cdots$), 试证明
\begin{align*}
m^*\left(\bigcup\limits_{n=1}^{\infty}B_n\right)=\sum\limits_{n=1}^{\infty}m^*(B_n).
\end{align*}
\end{proposition}
\begin{proof}
只需指出左端大于等于右端. 在公式
\begin{align*}
\sum\limits_{n=1}^{N}m^*(B_n)\xlongequal{\text{\refcor{corollary:有限个无交可测集与任意集合的并的外测度的性质}}}m^*\left(\bigcup\limits_{n=1}^{N}B_n\right)\leqslant m^*\left(\bigcup\limits_{n=1}^{\infty}B_n\right)
\end{align*}
中, 令 $N\to\infty$, 即得所证.
\end{proof}

\begin{theorem}[可测集的性质]\label{theorem:可测集的性质}
\begin{enumerate}[(1)]
\item \(\varnothing \in \mathscr{M}\).
\item 若 \(E \in \mathscr{M}\), 则 \(E^c \in \mathscr{M}\).

\item 若 \(E_1 \in \mathscr{M}\), \(E_2 \in \mathscr{M}\), 则 \(E_1 \cup E_2\), \(E_1 \cap E_2\) 以及 \(E_1 \setminus E_2\) 皆属于 \(\mathscr{M}\). (由此知, 可测集任何有限次取交、并运算后所得的集皆为可测集. )
\item 若 \(E_i \in \mathscr{M}\) (\(i = 1,2,\cdots\)), 则其并集$\bigcup_{i = 1}^{\infty} E_i$也属于 \(\mathscr{M}\). 若进一步有 \(E_i \cap E_j = \varnothing\) (\(i \neq j\)), 则
\begin{align*}
m^*\left(\bigcup_{i = 1}^{\infty} E_i\right) = \sum_{i = 1}^{\infty} m^*(E_i),
\end{align*}
即 \(m^*\) 在 \(\mathscr{M}\) 上满足\textbf{可数可加性}(或称为 \(\sigma -\)可加性).

\item\label{theorem:可测集的性质-5} 若 \(E_i \in \mathscr{M}\) (\(i = 1,2,\cdots\)), 则其交集$\bigcap_{i = 1}^{\infty} E_i$也属于 \(\mathscr{M}\). 

\item 如果$A$和$B$分别为$p$维和$q$维空间的可测集，那么$A\times B$是$p + q$维空间的可测集，测度为
\begin{align*}
m(A\times B)=m(A)\cdot m(B).
\end{align*}
\end{enumerate}
\end{theorem}
\begin{proof}
\begin{enumerate}[(1)]
\item 显然成立.
\item 注意到 \((E^c)^c = E\), 从定义可立即得出结论.
\item 对于任一集 \(T \subset \mathbb{R}^n\), 根据集合分解(参阅\hyperref[figure:集合关系示意图1]{图\ref{figure:集合关系示意图1}})与外测度的次可加性, 我们有
\begin{align*}
m^*(T) &\leqslant  m^*(T \cap (E_1 \cup E_2)) + m^*(T \cap (E_1 \cup E_2)^c) \\
&= m^*(T \cap (E_1 \cup E_2)) + m^*((T \cap E_1^c) \cap E_2^c) \\
&\leqslant  m^*((T \cap E_1) \cap E_2) + m^*((T \cap E_1) \cap E_2^c) \\
&\quad + m^*((T \cap E_1^c) \cap E_2) + m^*((T \cap E_1^c) \cap E_2^c).
\end{align*}
又由 \(E_1\), \(E_2\) 的可测性知, 上式右端就是
\begin{align*}
m^*(T \cap E_1) + m^*(T \cap E_1^c) = m^*(T).
\end{align*}
这说明
\begin{align*}
m^*(T) = m^*(T \cap (E_1 \cup E_2)) + m^*(T \cap (E_1 \cup E_2)^c).
\end{align*}
也就是说 \(E_1 \cup E_2\) 是可测集.
\begin{figure}[H]
\centering
\includegraphics[scale=0.3]{集合关系示意图1}
\caption{}
\label{figure:集合关系示意图1}
\end{figure}
为证 \(E_1 \cap E_2\) 是可测集, 只需注意 \(E_1 \cap E_2 = (E_1^c \cup E_2^c)^c\) 即可. 又由 \(E_1 \setminus E_2 = E_1 \cap E_2^c\) 可知, \(E_1 \setminus E_2\) 是可测集. 
\item 首先, 设 \(E_1\), \(E_2\), \(\cdots\), \(E_i\), \(\cdots\) 皆互不相交, 并令
\[
S = \bigcup_{i = 1}^{\infty} E_i, \quad S_k = \bigcup_{i = 1}^k E_i, \quad k = 1,2,\cdots.
\]
由(3)知每个 \(S_k\) 都是可测集, 从而对任一集 \(T\), 我们有
\begin{align*}
m^*(T) = m^*(T \cap S_k) + m^*(T \cap S_k^c) 
= m^*\left(\bigcup_{i = 1}^k (T \cap E_i)\right) + m^*(T \cap S_k^c) 
\xlongequal{\text{\refcor{corollary:有限个无交可测集与任意集合的并的外测度的性质}}} \sum_{i = 1}^k m^*(T \cap E_i) + m^*(T \cap S_k^c).
\end{align*}
由于 \(T \cap S_k^c \supset T \cap S^c\), 可知
\begin{align*}
m^*(T) \geqslant  \sum_{i = 1}^k m^*(T \cap E_i) + m^*(T \cap S^c).
\end{align*}
令 \(k \to \infty\), 就有
\begin{align*}
m^*(T) \geqslant  \sum_{i = 1}^{\infty} m^*(T \cap E_i) + m^*(T \cap S^c).
\end{align*}
再由外测度的次可加性可得
\begin{align*}
m^*(T)&\geqslant \sum_{i=1}^{\infty}{m^*(T}\cap E_i)+m^*(T\cap S^c)\geqslant m^*(\bigcup_{i=1}^{\infty}{\left( T\cap E_i \right)})+m^*(T\cap S^c)
\\
&=m^*(T\cap \bigcup_{i=1}^{\infty}{E_i})+m^*(T\cap S^c)=m^*(T\cap S)+m^*(T\cap S^c).
\end{align*}
这说明 \(S \in \mathscr{M}\).
此外, 在公式
\begin{align*}
m^*(T) \geqslant  \sum_{i = 1}^{\infty} m^*(T \cap E_i) + m^*(T \cap S^c)
\end{align*}
中以 \(T \cap S\) 替换 \(T\), 则又可得
\begin{align*}
m^*(T \cap S) \geqslant  \sum_{i = 1}^{\infty} m^*(T \cap E_i).
\end{align*}
又由外测度的次可加性可知反向不等式总是成立的, 因而实际上有
\begin{align*}
m^*(T \cap S) = \sum_{i = 1}^{\infty} m^*(T \cap E_i).
\end{align*}
在这里再取 \(T\) 为全空间 \(\mathbb{R}^n\), 就可证明可数可加性质:
\begin{align*}
m^*(S) = m^*\left(\bigcup_{i = 1}^{\infty} E_i\right) = \sum_{i = 1}^{\infty} m^*(E_i).
\end{align*}

其次, 对于一般的可测集列 \(\{E_i\}\), 我们令
\[
S_1 = E_1, \quad S_k = E_k \setminus \left(\bigcup_{i = 1}^{k - 1} E_i\right), \quad k \geqslant  2,
\]
则 \(\{S_k\}\) 是互不相交的可测集列. 而由 \(\bigcup_{i = 1}^{\infty} E_i = \bigcup_{k = 1}^{\infty} S_k\) 可知, \(\bigcup_{i = 1}^{\infty} E_i\) 是可测集.

\item 由(2)可知$E_i^c\in \mathscr{M}$,再由(4)可知$\bigcup_{i = 1}^{\infty} E_i^c$.于是再利用(2)和De Morgan定律可得
\begin{align*}
\left( \bigcup_{i=1}^{\infty}{E_{i}^{c}} \right) ^c=\bigcap_{i=1}^{\infty}{E_i}\in \mathscr{M} .
\end{align*}

\item 证明见\href{https://zhuanlan.zhihu.com/p/365546947#:~:text=%E5%AE%9A%E7%90%864.1%20%EF%BC%88%E5%8F%AF%E6%B5%8B%E9%9B%86%E7%9A%84%E7%9B%B4%E7%A7%AF%E4%BB%8D%E7%84%B6%E5%8F%AF%E6%B5%8B%EF%BC%8C%E5%B9%B6%E4%B8%94%E7%9B%B4%E7%A7%AF%E7%9A%84%E6%B5%8B%E5%BA%A6%E4%B8%BA%E6%B5%8B%E5%BA%A6%E7%9A%84%E4%B9%98%E7%A7%AF%EF%BC%89%20%E5%A6%82%E6%9E%9C%20A%20%E5%92%8C%20B%20%E5%88%86%E5%88%AB%E4%B8%BA%20p,%28A%29%20cdot%20m%20%28B%29%5C%20%E8%AF%81%E6%98%8E%E6%8C%89%E7%85%A7%E4%BB%A5%E4%B8%8B%E7%89%B9%E6%AE%8A%E5%88%B0%E4%B8%80%E8%88%AC%E7%9A%84%E8%BF%87%E7%A8%8B%E3%80%82%20%E6%96%B9%E4%BD%93%EF%BC%9A%E6%96%B9%E4%BD%93%E7%9A%84%E7%9B%B4%E7%A7%AF%E4%BB%8D%E7%84%B6%E6%98%AF%E6%96%B9%E4%BD%93%EF%BC%8C%E6%89%80%E4%BB%A5%E5%8F%AF%E6%B5%8B%E3%80%82%20%E5%85%B6%E6%B5%8B%E5%BA%A6%E5%B0%B1%E6%98%AF%E5%8E%9F%E6%9D%A5%E4%B8%A4%E4%B8%AA%E6%96%B9%E4%BD%93%E7%9A%84%E4%B9%98%E7%A7%AF%E3%80%82%20%E6%A0%B9%E6%8D%AE%E6%B5%8B%E5%BA%A6%E7%9A%84%E5%8F%AF%E5%8A%A0%E6%80%A7%E5%BE%97%E5%88%B0%E7%9B%B4%E7%A7%AF%E7%9A%84%E6%B5%8B%E5%BA%A6%E7%AD%89%E4%BA%8E%E6%B5%8B%E5%BA%A6%E7%9A%84%E4%B9%98%E7%A7%AF%E3%80%82}{知乎专栏}.
\end{enumerate} 
\end{proof}

\begin{corollary}\label{corollary:可测集类是sigma-代数}
$\mathscr{M}$是$\mathbb{R}^n$上的一个$\sigma$-代数.
\end{corollary}
\begin{proof}
由\hyperref[theorem:可测集的性质]{可测集的性质(1)(2)(4)}立得.
\end{proof}


\begin{proposition}\label{proposition:Cantor集可测且测度为0}
证明:Cantor集C是可测的,并且$m(C)=0$.
\end{proposition}
\begin{proof}
开区间是可测的. 由\hyperref[theorem:开集构造定理]{开集构造定理}, 我们知道 \(\mathbb{R}\) 中的开集是开区间的可数并, 因此也可测. 因此, 闭集也是可测的. 显然, 每个 \(C_n\) 都是闭集. 并且
\begin{align*}
C = \bigcap_{n = 1}^{\infty} C_n
\end{align*}
于是 \(C\) 也是闭集. 因此 \(C\) 是可测的.

下面, 我们用两种方法计算康托集的测度.

{\color{blue}法一:}根据我们的构造, \(C_{n + 1}\) 的测度刚好是去掉了 \(\frac{1}{3}\) 的 \(C_n\) 的测度. 换言之,
\begin{align*}
m(C_{n + 1}) = \left(1 - \frac{1}{3}\right)m(C_n) = \frac{2}{3}m(C_n)
\end{align*}
递归地, 对任意 \(n \in \mathbb{N}\), 我们有
\begin{align*}
m(C_n) = \left(\frac{2}{3}\right)^n m(C_0) = \left(\frac{2}{3}\right)^n
\end{align*}
注意到
\begin{align*}
m(C_0) = 1 < \infty
\end{align*}
因此由测度的第二单调收敛定理,
\begin{align*}
m(C) = m\left(\bigcap_{n = 1}^{\infty} C_n\right) = \lim_{n \to \infty} m(C_n) = \lim_{n \to \infty} \left(\frac{2}{3}\right)^n = 0
\end{align*}
此即得证.

{\color{blue}法二:}设 \(n \geqslant  2\). \(C_n\) 比 \(C_{n - 1}\) 减少了 \(2^{n - 1}\) 个区间, 每个区间长度为 \(\frac{1}{3^n}\). 因此 \(C_n\) 比 \(C_{n - 1}\) 减少的长度为
\begin{align*}
2^{n - 1}\frac{1}{3^n} = \frac{1}{3}\left(\frac{2}{3}\right)^{n - 1}
\end{align*}
总共减少的长度为
\begin{align*}
\sum_{n = 1}^{\infty} \frac{1}{3}\left(\frac{2}{3}\right)^{n - 1} = \frac{1}{3} \frac{1}{1 - \frac{2}{3}} = \frac{1}{3} \cdot 3 = 1
\end{align*}
因此
\begin{align*}
m(C) = 1 - 1 = 0.
\end{align*} 
\end{proof}

\begin{proposition}\label{proposition:可测集类的基数是2^c}
\(\mathscr{M}\) 的基数是 \(2^c\).
\end{proposition}
\begin{proof}
由\refpro{proposition:Cantor集可测且测度为0}可知Cantor集是零测集,不难推断 \(\mathscr{M}\) 的基数大于或等于 \(2^c\), 但 \(\mathscr{M}\) 的基数又不会超过 \(2^c\), 于是 \(\mathscr{M}\) 的基数实际上是 \(2^c\). 
\end{proof}

\begin{definition}[Lebesgue测度]
对于可测集 \(E\), 其外测度称为\textbf{测度}, 记为 \(m(E)\). 这就是通常所说的 \(\mathbb{R}^n\) 上的 \textbf{Lebesgue测度}.
\end{definition}

\begin{definition}[测度]
设 \(X\) 是非空集合, \(\mathscr{A}\) 是 \(X\) 的一些子集构成的 \(\sigma -\)代数. 若 \(\mu\) 是定义在 \(\mathscr{A}\) 上的一个集合函数, 且满足：
\begin{enumerate}[(i)]
\item \(0 \leqslant  \mu(E) \leqslant  +\infty\) (\(E \in \mathscr{A}\));
\item \(\mu(\varnothing)=0\);
\item \(\mu\) 在 \(\mathscr{A}\) 上是可数可加的,
\end{enumerate}
则称 \(\mu\) 是 \(\mathscr{A}\) 上的(非负)\textbf{测度}. \(\mathscr{A}\) 中的元素称为(\(\mu\))\textbf{可测集}, 有序组 \((X, \mathscr{A}, \mu)\) 称为\textbf{测度空间}. 
\end{definition}
\begin{remark}
由\refcor{corollary:可测集类是sigma-代数}可知$\mathscr{M}$就是$\mathbb{R}^n$上的一个$\sigma$-代数,故本节所建立的测度空间就是 \((\mathbb{R}^n, \mathscr{M}, m)\). 
\end{remark}

\begin{theorem}[测度的基本性质]\label{theorem:测度的基本性质}
\begin{enumerate}[(1)]
\item 非负性: 若$E\in \mathscr{M}$,则\(m(E) \geqslant  0\), \(m(\varnothing)=0\);

\item\label{theorem:测度的基本性质-2} 单调性: 若$E_1,E_2\in \mathscr{M}$且\(E_1 \subset E_2\), 则 \(m(E_2) \leqslant  m(E_1)\).

若只有$E_1\in \mathscr{M}$,则$m^*\left( E_2 \right) =m^*\left( E_1 \right) +m^*\left( E_2\setminus E_1 \right)$.

并且若还有$m(E_1)<+\infty$,则$m^*(E_2\setminus E_1)+m^*(E_1)=m^*(E_2)$.

\item 可数可加性: 若 \(E_i \in \mathscr{M}\) (\(i = 1,2,\cdots\))且 \(E_i \cap E_j = \varnothing\) (\(i \neq j\)),则
\begin{align*}
m\left(\bigcup_{i = 1}^{\infty} E_i\right) = \sum_{i = 1}^{\infty} m(E_i).
\end{align*}

\item\label{theorem:测度的基本性质-4} 若$E_1,E_2\in \mathscr{M}$,且$m(E_1\cap E_2)<+\infty$,则$m(E_1\cup E_2)=m(E_1)+m(E_2)-m(E_1\cap E_2)$.
\end{enumerate}
\end{theorem}
\begin{proof}
\begin{enumerate}[(1)]
\item 由\hyperref[theorem:R^n 中点集的外测度性质]{$R^n$ 中点集的外测度性质}立得.

\item 由\hyperref[theorem:R^n 中点集的外测度性质]{$R^n$ 中点集的外测度性质}可知\(m(E_1) \leqslant  m(E_2)\).再根据$E_1$可测的定义可知
\begin{align*}
m^*\left( E_2 \right) =m^*\left( E_2\cap E_1 \right) +m^*\left( E_2\cap E_{1}^{c} \right) =m^*\left( E_1 \right) +m^*\left( E_2\setminus E_1 \right) .
\end{align*}
又$m(E_1)<+\infty$,故由上式移项后得$m(E_2\backslash E_1)=m(E_2)-m(E_1).$

\item 由\hyperref[theorem:可测集的性质]{可测集的性质}立得.

\item 注意到$E_1\cap (E_2\backslash (E_1\cap E_2))=\varnothing$,并且$E_1\cap E_2\subset E_2,m(E_1\cap E_2)<+\infty$,故由(2)(3)可得
\begin{align*}
m\left( E_1\cup E_2 \right) =m\left[ E_1\cup \left( E_2\backslash (E_1\cap E_2) \right) \right] =m\left( E_1 \right) +m\left[ E_2\backslash (E_1\cap E_2) \right] =m\left( E_1 \right) +m\left( E_2 \right) -m\left( E_1\cap E_2 \right) .
\end{align*}
\end{enumerate}
\end{proof}

\begin{theorem}[递增可测集列的测度运算]\label{theorem:递增可测集列的测度运算}
设 $(X,\mathscr{M},\mu)$ 为测度空间,若有递增可测集列$E_1\subset E_2\subset\cdots\subset E_k\cdots\subset \mathscr{M}$，则
\begin{align}
m\left(\lim_{k\to\infty}E_k\right)=\lim_{k\to\infty}m(E_k).\label{eq:2.4}
\end{align}
\end{theorem}
\begin{proof}
若存在$k_0$，使得$m(E_{k_0})=+\infty$，则
\begin{align*}
m^*\left(\lim_{k\rightarrow \infty}E_k\right) = m^*\left(\bigcup_{k=1}^{\infty}E_k\right)\geqslant m^*(E_{k_0}).
\end{align*}
因此$m^*\left(\lim_{k\rightarrow \infty}E_k\right) = +\infty$. 又由$\{E_k\}_{k=1}^{\infty}$递增可知
\begin{align*}
m^*(E_k) \geqslant m^*(E_{k_0}), \quad \forall k \geqslant k_0.
\end{align*}
因此$\lim_{k\rightarrow \infty}m^*(E_k) = +\infty$.
故此时定理自然成立。

现在假定对一切$k$，有$m(E_k)<+\infty$。由假设$E_k\in\mathscr{M}$（$k = 1,2,\cdots$），故$E_{k - 1}$与$E_k\setminus E_{k - 1}$是互不相交的可测集。由测度的可加性知
$m(E_{k - 1})+m(E_k\setminus E_{k - 1})=m(E_k)$。
因为$m(E_{k - 1})$是有限的，所以移项得
$m(E_k\setminus E_{k - 1})=m(E_k)-m(E_{k - 1})$。
令$E_0 = \varnothing$，可得
$\lim_{k\to\infty}E_k=\bigcup_{k = 1}^{\infty}E_k=\bigcup_{k = 1}^{\infty}(E_k\setminus E_{k - 1})$。
再应用\hyperref[theorem:测度的基本性质]{测度的可数可加性}，我们有
\begin{align*}
m\left(\lim_{k\to\infty}E_k\right)&=m\left(\bigcup_{k = 1}^{\infty}(E_k\setminus E_{k - 1})\right)
=\sum_{k = 1}^{\infty}(m(E_k)-m(E_{k - 1}))\\
&=\lim_{k\to\infty}\sum_{i = 1}^{k}(m(E_i)-m(E_{i - 1}))
=\lim_{k\to\infty}m(E_k).
\end{align*}
\end{proof}

\begin{corollary}[递减可测集列的测度运算]\label{corollary:递减可测集列的测度运算}
若有递减可测集列$E_1\supset E_2\supset\cdots\supset E_k\supset\cdots$，且$m(E_1)<+\infty$，则
\begin{align}
m\left(\lim_{k\to\infty}E_k\right)=\lim_{k\to\infty}m(E_k).\label{eq:2.5}
\end{align}
\end{corollary}
\begin{proof}
由\hyperref[theorem:可测集的性质]{可测集的性质(5)}可知$\lim_{k\to\infty}E_k$是可测集,再由\hyperref[theorem:测度的基本性质]{测度的单调性}可知$\lim_{k\to\infty}m(E_k)\leqslant  m(E_1)<+\infty$.因为
$E_1\setminus E_k\subset E_1\setminus E_{k + 1},k = 2,3,\cdots$，
所以由\hyperref[theorem:可测集的性质]{可测集的性质(2)}可知$\{E_1\setminus E_k\}$是递增可测集合列。于是由\hyperref[theorem:递增可测集列的测度运算]{递增可测集列的测度运算}可知
\begin{align*}
m\left(E_1\setminus\lim_{k\to\infty}E_k\right)=m\left(\lim_{k\to\infty}(E_1\setminus E_k)\right)
=\lim_{k\to\infty}m(E_1\setminus E_k).
\end{align*}
由于$m(E_1)<+\infty$，故由\hyperref[theorem:测度的基本性质]{测度的基本性质(2)}上式可写为
$m(E_1)-m\left(\lim_{k\to\infty}E_k\right)=m(E_1)-\lim_{k\to\infty}m(E_k)$。
消去$m(E_1)$，我们有
$m\left(\lim_{k\to\infty}E_k\right)=\lim_{k\to\infty}m(E_k)$。 
\end{proof}

\begin{theorem}[Borel-Cantelli引理]\label{theorem:Borel-Cantelli引理}
\begin{enumerate}[(1)]
\item\label{theorem:Borel-Cantelli引理-1} 若有可测集列$\{E_k\}$，且有$\sum_{k = 1}^{\infty}m(E_k)<+\infty$，则
\begin{align*}
m\left(\varlimsup_{k\to\infty}E_k\right)= 0.
\end{align*}

\item\label{theorem:Borel-Cantelli引理-2} 设$\{E_k\}$是可测集列，则
\begin{align*}
m\left(\varliminf_{k\to\infty}E_k\right)\leqslant\varliminf_{k\to\infty}m(E_k).
\end{align*}
当存在$n\in \mathbb{N}$,使得$m\left( \bigcup_{k=n}^{\infty}{E_k} \right) <+\infty$时,有
\begin{align*}
m\left(\varlimsup_{k\to\infty}E_k\right)\geqslant\varlimsup_{k\to\infty}m(E_k).
\end{align*}
\end{enumerate}
\end{theorem}
\begin{proof}
\begin{enumerate}[(1)]
\item 由\hyperref[corollary:递减可测集列的测度运算]{递减可测集列的测度运算}可得
\begin{align*}
m\left( \underset{k\rightarrow \infty}{\overline{\lim }}E_k \right) &=m\left( \bigcap_{k=1}^{\infty}{\bigcup_{i=k}^{\infty}{E_i}} \right) =m\left( \lim_{k\rightarrow \infty} \bigcup_{i=k}^{\infty}{E_i} \right) 
\\
&\xlongequal{\text{\hyperref[corollary:递减可测集列的测度运算]{递减可测集列的测度运算}}}\lim_{k\rightarrow \infty} m\left( \bigcup_{i=k}^{\infty}{E_i} \right) \\
&\leqslant \lim_{k\rightarrow \infty} \sum_{i=k}^{\infty}{m(E_i)}=0.
\end{align*}

\item 因为$\bigcap_{j = k}^{\infty}E_j\subset E_k,\bigcup_{i = k}^{\infty}E_i\supset E_k$（$k = 1,2,\cdots$），所以有
\begin{align*}
m\left(\bigcap_{j = k}^{\infty}E_j\right)\leqslant m(E_k),\quad m\left(\bigcup_{i = k}^{\infty}E_i\right)\geqslant m(E_k)\quad (k = 1,2,\cdots).
\end{align*}
令$k\to\infty$，则得（$\bigcap_{j = k}^{\infty}E_j$随$k$增大而递增,$\bigcup_{i = k}^{\infty}E_i$随$k$增大而递减）
\begin{align*}
m\left( \underset{k\rightarrow \infty}{\underline{\lim }}E_k \right) &=m\left( \bigcup_{k=1}^{\infty}{\bigcap_{j=k}^{\infty}{E_j}} \right) =m\left( \lim_{k\rightarrow \infty} \bigcap_{j=k}^{\infty}{E_j} \right) 
\\
&\xlongequal{\text{\hyperref[theorem:递增可测集列的测度运算]{递增可测集列的测度运算}}}\lim_{k\rightarrow \infty} m\left( \bigcap_{j=k}^{\infty}{E_j} \right) \\
&\leqslant \underset{k\rightarrow \infty}{\underline{\lim }}m(E_k).
\end{align*}
\begin{align*}
m\left( \underset{k\rightarrow \infty}{\overline{\lim }}E_k \right) &=m\left( \bigcap_{k=1}^{\infty}{\bigcup_{i=k}^{\infty}{E_i}} \right) =m\left( \lim_{k\rightarrow \infty} \bigcup_{i=k}^{\infty}{E_i} \right) 
\\
&\xlongequal{\text{\hyperref[corollary:递减可测集列的测度运算]{递减可测集列的测度运算}}}\lim_{k\rightarrow \infty} m\left( \bigcup_{i=k}^{\infty}{E_i} \right) \\
&\geqslant \underset{k\rightarrow \infty}{\overline{\lim }}m(E_k).
\end{align*}
\end{enumerate} 
\end{proof}

\begin{example}
设有 $\mathbb{R}$ 中可测集列 $\{E_k\}$, 且当 $k\geqslant k_0$ 时, $E_k\subseteq[a,b]$. 若存在 $\lim\limits_{k\to\infty}E_k=E$, 试证明: $m(E)=\lim\limits_{k\to\infty}m(E_k)$.
\end{example}
\begin{proof}
由条件知$\bigcup\limits_{k=k_0}^{\infty}E_k\subseteq \left[ a,b \right]$,从而$m\left( \bigcup\limits_{k=k_0}^{\infty}E_k \right) \leqslant b-a<+\infty$.于是由\hyperref[theorem:Borel-Cantelli引理]{Borel-Cantelli引理}知
\begin{align*}
m\left( E \right) =m\left( \underset{k\rightarrow \infty}{\overline{\lim }}E_k \right) \geqslant \underset{k\rightarrow \infty}{\overline{\lim }}m\left( E_k \right) \geqslant \underset{k\rightarrow \infty}{\underline{\lim }}m\left( E_k \right) \geqslant m\left( \underset{k\rightarrow \infty}{\underline{\lim }}E_k \right) =m\left( E \right) .
\end{align*}
故$m\left( E \right) =\lim\limits_{k\rightarrow \infty}m\left( E_k \right).$
\end{proof}

\begin{proposition}\label{proposition:可测集的一些性质111}
\begin{enumerate}[(1)]
\item 设 $E\subseteq\mathbb{R}$, $0\neq\lambda\in\mathbb{R}$. 记 $\lambda E=\{\lambda x: x\in E\}$. 若 $E\in\mathscr{M} $, 则 $\lambda E\in\mathscr{M} $.

\item\label{proposition:可测集的一些性质111-2} 设 $E\subseteq\mathbb{R}^n$. 若存在 $\mathbb{R}^n$ 中可测集 $A$, 使得 $m^*(E\triangle A)=0$, 则 $E\in\mathscr{M} $, 且有 $m(E)=m(A)$.

\item 设 $A_1,A_2\subseteq\mathbb{R}^n$, $A_1\subseteq A_2$, $A_1$ 是可测集且有 $m(A_1)=m^*(A_2)<+\infty$, 则 $A_2$ 是可测集.

\item 设 $E\subseteq\mathbb{R}^n$. 若对任给 $\varepsilon>0$, 均存在 $A\in\mathscr{M} $, 使得 $m^*(E\triangle A)<\varepsilon$, 则 $E\in\mathscr{M} $.
\end{enumerate}
\end{proposition}
\begin{proof}
\begin{enumerate}[(1)]
\item 因为对任意的试验集 $T\subseteq\mathbb{R}$, 有
\begin{align*}
T\cap\lambda E=\lambda(\lambda^{-1}T\cap E),\quad T\cap(\lambda E)^c=\lambda(\lambda^{-1}T\cap E^c),
\end{align*}
所以可得
\begin{align*}
m^*(T\cap\lambda E)+m^*(T\cap(\lambda E)^c)\xlongequal{\text{\hyperref[theorem:外测度的数乘]{外测度的数乘}}}|\lambda|\,[m^*(\lambda^{-1}T\cap E)+m^*(\lambda^{-1}T\cap E^c)].
\end{align*}
这说明 $E$ 的可测性等价于 $\lambda E$ 的可测性.

\item 注意到
\begin{align*}
E\setminus A,A\setminus E\subseteq (E\setminus A)\cup (A\setminus E)=E\triangle A,
\end{align*}
故
\begin{align*}
m^*(E\setminus A),m^*(A\setminus E)\leqslant m^*(E\triangle A)=0.
\end{align*}
从而$E\setminus A,A\setminus E$都是零测集，进而都可测.
又注意到$E=(E\setminus A)\cup [A\setminus (A\setminus E)]$，故$E$也可测，且
\begin{align*}
m(E)=m(E\setminus A)+m(A)-m(A\setminus E)=m(A).
\end{align*}

\item 由题设知
\begin{align*}
m(A_1)=m^*(A_2)=m^*(A_1\cup(A_2\setminus A_1))
=m^*(A_1)+m^*(A_2\setminus A_1),
\end{align*}
故可得 $m^*(A_2\setminus A_1)=0$. 从而 $A_2\setminus A_1$ 是可测集, 即 $A_2=A_1\cup(A_2\setminus A_1)$ 是可测集.

\item 依题设知, 对任给 $\varepsilon>0$, 取 $\frac{\varepsilon}{2^k}$ ($k\in\mathbb{N}$), 存在可测集列 $\{A_k\}$, 使得 $m^*(E\triangle A_k)<\frac{\varepsilon}{2^k}$ ($k\in\mathbb{N}$), 即有 $m^*(E\setminus A_k)<\frac{\varepsilon}{2^k}$, $m^*(A_k\setminus E)<\frac{\varepsilon}{2^k}$ ($k\in\mathbb{N}$). 现在令 $A=\bigcup\limits_{k=1}^{\infty}A_k$, 则 $A$ 是可测集, 且有
\begin{align*}
E\setminus A=\bigcap\limits_{k=1}^{\infty}(E\setminus A_k),\quad A\setminus E=\bigcup\limits_{k=1}^{\infty}(A_k\setminus E),
\end{align*}
\begin{align*}
m^*(E\setminus A)\leqslant m^*(E\setminus A_k)\ (k\in\mathbb{N}),\quad m^*(A\setminus E)\leqslant\sum\limits_{k=1}^{\infty}m^*(A_k\setminus E),
\end{align*}
\begin{align*}
m^*(E\setminus A)\leqslant\frac{\varepsilon}{2^k}\ (k\in\mathbb{N}),\quad m^*(A\setminus E)\leqslant\sum\limits_{k=1}^{\infty}\frac{\varepsilon}{2^k}=\varepsilon.
\end{align*}
由 $\varepsilon$ 的任意性, 可知 $m^*(E\setminus A)=0=m^*(A\setminus E)$, 即 $m^*(E\triangle A)=0$. 根据\rrefpro{proposition:可测集的一些性质111}{proposition:可测集的一些性质111-2}知,$E$ 是可测集.
\end{enumerate}
\end{proof}

\begin{proposition}
\begin{enumerate}[(1)]
\item 设 $A\in\mathscr{M}$,则对任意的 $B\subseteq\mathbb{R}^n$,必有
\begin{align*}
m^*(A\cup B)+m^*(A\cap B)=m^*(A)+m^*(B).
\end{align*}
\item 设 $A,B,C$ 是 $\mathbb{R}^n$ 中的可测集.若有 $m(A\triangle B)=0$, $m(B\triangle C)=0$,则 $m(A\triangle C)=0$.
\item 设 $E\subseteq[0,1]$.若 $m(E)=1$,试证明 $\overline{E}=[0,1]$;若 $m(E)=0$,则 $\mathring{E}=\varnothing$.
\item 设 $E_1,E_2,\cdots,E_k$ 是 $[0,1]$ 中的可测集,且有 $\sum\limits_{i=1}^{k}m(E_i)>k-1$,则 $m\left(\bigcap\limits_{i=1}^{k}E_i\right)>0$.
\end{enumerate}
\end{proposition}
\begin{proof}
\begin{enumerate}[(1)]
\item 当$m^*(A)=+\infty$时,有
\begin{align*}
m^*(A\cup B)\geqslant m^*(A)= +\infty. 
\end{align*}
故此时结论成立.
下设$m^*(A)<+\infty$,因为$A$可测,所以我们有
\begin{gather*}
m^*(B)=m^*(B\cap A)+m^*(B\cap A^c),
\\
m^*(B\cup A)=m^*((B\cup A)\cap A)+m^*((B\cup A)\cap A^c)
=m^*(A)+m^*(B\cap A^c).
\end{gather*}
由此可知 $m^*(B\cap A^c)=m^*(B\cup A)-m^*(A)$,从而得
\begin{align*}
m^*(B)=m^*(B\cap A)+m^*(B\cup A)-m^*(A).
\end{align*}
移项后即得所证.

\item 注意
\begin{gather*}
A\backslash C=A\cap C^c\cap \left( B\cup B^c \right) =\left[ A\cap \left( B\backslash C \right) \right] \cup \left[ C^c\cap \left( A\backslash B \right) \right] ,
\\
C\backslash A=C\cap A^c\cap \left( B\cup B^c \right) =\left[ C\cap \left( B\backslash A \right) \right] \cup \left[ A^c\cap \left( C\backslash B \right) \right] .
\end{gather*}
由$m(A\triangle B)=m(B\triangle C)=0$以及
\begin{align*}
A\setminus B,\,B\setminus A\subseteq A\triangle B,\quad B\setminus C,\,C\setminus B\subseteq B\triangle C,
\end{align*}
知
\begin{align*}
m(A\setminus B),\,m(B\setminus A)\leqslant m(A\triangle B)=0,\\
m(B\setminus C),\,m(C\setminus B)\leqslant m(B\triangle C)=0.
\end{align*}
又$A\triangle C=(A\setminus C)\cup (C\setminus A)$，于是
\begin{align*}
&m(A\triangle C)\leqslant m(A\setminus C)+m(C\setminus A)\\
&\leqslant m(A\cap (B\setminus C))+m(C^c\cap (A\setminus B))+m(C\cap (B\setminus A))+m(A^c\cap (C\setminus B))\\
&\leqslant m(B\setminus C)+m(A\setminus B)+m(B\setminus A)+m(C\setminus B)=0.
\end{align*}

\item 反证.
若$x_0\in [0,1]\setminus \overline{E}$,则由闭包的定义知,存在$\delta>0$,使得$B(x_0,\delta)\cap E=\varnothing$且$B(x_0,\delta)\subseteq [0,1]$.于是$B(x_0,\delta)\subseteq [0,1]\setminus E$,从而$0<2\delta=m(B(x_0,\delta))\leqslant m([0,1]\setminus E)=0$,矛盾!

若$\mathring{E}\neq \varnothing$,设$x_0\in \mathring{E}$,则存在$\delta>0$,使得$B(x_0,\delta)\subseteq E$.于是$0<2\delta=m(B(x_0,\delta))\leqslant m(E)=0$,矛盾!

\item 令 $A_i=[0,1]\setminus E_i$ ($i=1,2,\cdots,k$),则可得
\begin{align*}
\bigcap\limits_{i=1}^{k}E_i=[0,1]\setminus\bigcup\limits_{i=1}^{k}A_i,\quad m\left(\bigcap\limits_{i=1}^{k}E_i\right)=1-m\left(\bigcup\limits_{i=1}^{k}A_i\right).
\end{align*}
注意到
\begin{align*}
m\left(\bigcup\limits_{i=1}^{k}A_i\right)\leqslant\sum\limits_{i=1}^{k}m(A_i)=\sum\limits_{i=1}^{k}(1-m(E_i))
=k-\sum\limits_{i=1}^{k}m(E_i)<k-(k-1)=1,
\end{align*}
即知 $m\left(\bigcap\limits_{i=1}^{k}E_i\right)>0$.
\end{enumerate}
\end{proof}

\begin{example}
设 $\{E_n\}$ 是 $[0,1]$ 中互不相同的可测集合列, 且存在 $\varepsilon\in (0,1)$, $m(E_n)\geqslant\varepsilon$ ($n=1,2,\cdots$). 试问是否存在子列 $\{E_{n_i}\}$, 使得
\begin{align*}
m\left(\bigcap\limits_{i=1}^{\infty}E_{n_i}\right)>0?
\end{align*}
\end{example}
\begin{proof}
不一定. 例如作点集列: $E_n$ ($n\in\mathbb{N}$),
\[
E_n=\{x\in[0,1]: x\text{ 的十进位小数表示式中, 第 }n\text{ 位数字 }\neq 0\},
\]
从而对$\forall k\in \mathbb{N}$,都有
\begin{align*}
\bigcap\limits_{i=1}^k{E_{n_i}}=\left\{ x\in [0,1]:x\text{的十进制小数表示式中,第}n_i\left( i=1,2,\cdots ,k \right) \text{位数字都不为}0 \right\}.
\end{align*}
取$N=\max\left\{ n_1,n_2,\cdots ,n_k \right\}$,对$[0,1]$区间作$10^N$等分,并记
\begin{align*}
I_{d_1d_2\cdots d_N}=\left[ \frac{d_1}{10}+\frac{d_2}{10}+\cdots +\frac{d_N}{10^N},\frac{d_1}{10}+\frac{d_2}{10}+\cdots +\frac{d_N}{10^N}+\frac{1}{10^N} \right),\text{其中}d_i\in \{0,1,\cdots ,9\}.
\end{align*}
显然上述区间两两无交,且$m\left( I_{d_1d_2\cdots d_N} \right) =\frac{1}{10^N}$,$\forall d_i\in \{0,1,\cdots ,9\}$.于是
\begin{align*}
\bigcap\limits_{i=1}^k{E_{n_i}}=\bigcup\limits_{d_{n_1},d_{n_2},\cdots ,d_{n_k}\ne 0}{I_{d_1d_2\cdots d_N}}=\bigcup\limits_{d_{n_1},d_{n_2},\cdots ,d_{n_k}\in \{1,2,\cdots ,9\}}{I_{d_1d_2\cdots d_N}}.
\end{align*}
故
\begin{align*}
m\left( \bigcap_{i=1}^k{E_{n_i}} \right) 
&= m\left( \bigcup_{\substack{d_{n_1},\cdots,d_{n_k}\in\{1,\cdots,9\} \\ d_i\in\{0,\cdots,9\},\,i\neq n_1,\cdots,n_k}} I_{d_1d_2\cdots d_N} \right) 
= \sum_{\substack{d_{n_1},\cdots,d_{n_k}\in\{1,\cdots,9\} \\ d_i\in\{0,\cdots,9\},\,i\neq n_1,\cdots,n_k}} m\left( I_{d_1d_2\cdots d_N} \right) \\
&= \sum_{\substack{d_{n_1},\cdots,d_{n_k}\in\{1,\cdots,9\} \\ d_i\in\{0,\cdots,9\},\,i\neq n_1,\cdots,n_k}} \frac{1}{10^N} 
= 9^k \cdot 10^{N-k} \cdot \frac{1}{10^N} = \left( \frac{9}{10} \right)^k.
\end{align*}
令$k\to +\infty$,得
\begin{align*}
m\left( \bigcap_{i=1}^{\infty}{E_{n_i}} \right)=0. 
\end{align*}
\end{proof}

\begin{example}
设 $\{E_n\}$ 是 $[0,1]$ 中的可测集列, 且满足 $\varlimsup_{n\to\infty}m(E_n)=1$, 试证明对任意的 $\alpha$: $0<\alpha<1$, 必存在 $\{E_{n_k}\}$, 使得 $m\left(\bigcap\limits_{k=1}^{\infty}E_{n_k}\right)>\alpha$.
\end{example}
\begin{proof}
由题设知, 对任意的 $k\in\mathbb{N}$, 存在 $\{n_k\}$, 使得
\begin{align*}
m(E_{n_k})>1-\frac{1-\alpha}{2^k}\quad (k\in\mathbb{N}).
\end{align*}
由此知 $1-m(E_{n_k})<\frac{1-\alpha}{2^k}$ ($k\in\mathbb{N}$). 从而得到
\begin{align*}
[0,1]\setminus\bigcap\limits_{k=1}^{\infty}E_{n_k}=\bigcup\limits_{k=1}^{\infty}([0,1]\setminus E_{n_k}),
\end{align*}
\begin{align*}
m\left([0,1]\setminus\bigcap\limits_{k=1}^{\infty}E_{n_k}\right)\leqslant\sum\limits_{k=1}^{\infty}m([0,1]\setminus E_{n_k})
=\sum\limits_{k=1}^{\infty}(1-m(E_{n_k}))\leqslant\sum\limits_{k=1}^{\infty}\frac{1-\alpha}{2^k}=1-\alpha,
\end{align*}
故有
\begin{align*}
m\left(\bigcap\limits_{k=1}^{\infty}E_{n_k}\right)>\alpha.
\end{align*}
\end{proof}

\begin{example}
解答下列命题：
\begin{enumerate}[(1)]
\item 设$\{E_k\}$是$[0,1]$中的可测集列，$m(E_k)=1$（$k=1,2,\cdots$），试证明$m\left(\bigcap\limits_{k=1}^{\infty}E_k\right)=1$.

\item 试作$[0,1]$中的第二纲集$E$：$m(E)=0$.
\end{enumerate}
\end{example}
\begin{proof}
\begin{enumerate}[(1)]
\item 由题设知$m([0,1]\setminus E_k)=0$（$k\in\mathbb{N}$），又有
\begin{align*}
m\left([0,1]\setminus\bigcap\limits_{k=1}^{\infty}E_k\right) &= m\left(\bigcup\limits_{k=1}^{\infty}\bigl([0,1]\setminus E_k\bigr)\right) \leqslant \sum\limits_{k=1}^{\infty}m([0,1]\setminus E_k) = 0.
\end{align*}
由此得$m\left(\bigcap\limits_{k=1}^{\infty}E_k\right)=1$.

\item 令$[0,1]\cap\mathbb{Q}=\{r_1,r_2,\cdots,r_n,\cdots\}$，以及
\[
I_{n,k} = \left(r_n - \frac{1}{2^{n+k}},\, r_n + \frac{1}{2^{n+k}}\right) \quad (n,k\in\mathbb{N}),
\]
则点集$(-\infty,+\infty)\setminus\bigcup\limits_{n,k=1}^{\infty}I_{n,k}$在$\mathbb{R}$中无处稠密。我们有
\begin{align*}
m\left(\bigcap\limits_{k=1}^{\infty}\bigcup\limits_{n=1}^{\infty}I_{n,k}\right) = 0,
\end{align*}
$\bigcap\limits_{k=1}^{\infty}\bigcup\limits_{n=1}^{\infty}I_{n,k}$是第二纲集。
\end{enumerate}
\end{proof}

\begin{example}
解答下列问题：
\begin{enumerate}[(1)]
\item 将$(0,1)$中的数用十进位小数展开，求下列点集之测度。
\begin{enumerate}[(i)]
\item 在指定小数位置上是数字$4$的点之全体$E_1$.
\item 在指定两个小数位置上都是已给定的数字之全体$E_2$.
\item 在两个指定小数位置上是不同的数字之全体$E_3$.
\end{enumerate}

\item 将$(0,1)$中的数用十进位小数展开，在表示式中总有数字''$0$''出现的数的全体为$E$，求$m(E)$.
\item 在$[0,1]$中作点集
\[
E=\{x\in[0,1]: \text{在十进位小数表示式 }x=0.a_1a_2\cdots\text{中的所有 }a_i\text{都不出现 }10\text{个数字中的某一个}\},
\]
试证明$E$是不可数集，且$m(E)=0$.
\item 将$(0,1)$中的点用十进位小数展开，令表示式中不超过$9$个数字的点的全体为$E$，求$\overline{E}$以及$m(E)$.
\item 将$[0,1]$中的点用十进位小数展开，令
\[
E = \{x\in[0,1]: x\text{的任一位小数是 }2\text{ 或 }7\},
\]
试问：
\begin{enumerate}[(i)]
\item $E$是闭集还是开集？

\item $E$是可数集？

\item $\overline{E}=[0,1]$？

\item $E$是可测集？$m(E)=$？
\end{enumerate}
\end{enumerate}
\end{example}
\begin{solution}
\begin{enumerate}[(1)]
\item 不妨设
\begin{gather*}
E_1=\left\{ x\in \left[ 0,1 \right] :x\text{的十进位小数表示式中},\text{第}n_1\text{位数字是}4 \right\} ;
\\
E_2=\left\{ x\in \left[ 0,1 \right] :x\text{的十进位小数表示式中},\text{第}n_2,n_3\text{位数字是}m \right\} ,\text{其中}m\in \left\{ 0,1,2\cdots ,9 \right\} ;
\\
E_3=\left\{ x\in \left[ 0,1 \right] :x\text{的十进位小数表示式中},\text{第}n_4,n_5\text{位数字不同} \right\} .
\end{gather*}
取$N=\max\left\{ n_1,n_2,n_3,n_4,n_5 \right\}$,对$[0,1]$区间作$10^N$等分,并记
\begin{align*}
I_{d_1d_2\cdots d_N}=\left[ \frac{d_1}{10}+\frac{d_2}{10}+\cdots +\frac{d_N}{10^N},\frac{d_1}{10}+\frac{d_2}{10}+\cdots +\frac{d_N}{10^N}+\frac{1}{10^N} \right),\text{其中}d_i\in \{0,1,\cdots ,9\}.
\end{align*}
显然上述区间两两无交,且$m\left( I_{d_1d_2\cdots d_N} \right) =\frac{1}{10^N}$,$\forall d_i\in \{0,1,\cdots ,9\}$.于是
\begin{align*}
m\left( E_1 \right) &=m\left( \bigcup\limits_{\substack{d_{n_1}=4,\\d_i\in \left\{ 0,1,\cdots ,9 \right\} ,i\ne n_1}}I_{d_1d_2\cdots d_N} \right) =\sum\limits_{\substack{d_{n_1}=4,\\d_i\in \left\{ 0,1,\cdots ,9 \right\} ,i\ne n_1}}m\left( I_{d_1d_2\cdots d_N} \right)=10^{N-1}\cdot \frac{1}{10^N}=\frac{1}{10},\\
m\left( E_2 \right) &=m\left( \bigcup\limits_{\substack{d_{n_2}=d_{n_3}=4,\\d_i\in \left\{ 0,1,\cdots ,9 \right\} ,i\ne n_2,n_3}}I_{d_1d_2\cdots d_N} \right) =\sum\limits_{\substack{d_{n_2}=d_{n_3}=4,\\d_i\in \left\{ 0,1,\cdots ,9 \right\} ,i\ne n_2,n_3}}m\left( I_{d_1d_2\cdots d_N} \right)=10^{N-2}\cdot \frac{1}{10^N}=\frac{1}{100},\\
m\left( E_3 \right) &=m\left( \bigcup\limits_{\substack{d_{n_4}\ne d_{n_5},\\d_i\in \left\{ 0,1,\cdots ,9 \right\} ,i\ne n_4,n_5}}I_{d_1d_2\cdots d_N} \right) =\sum\limits_{\substack{d_{n_4}\ne d_{n_5},\\d_i\in \left\{ 0,1,\cdots ,9 \right\} ,i\ne n_4,n_5}}m\left( I_{d_1d_2\cdots d_N} \right)=\left( 10^2-10 \right) \cdot 10^{N-2}\cdot \frac{1}{10^N}=\frac{9}{10}.
\end{align*}

\item 对任意的$n\in\mathbb{N}$，作点集
\[
E_n = \{x\in(0,1): x\text{的小数展开式中前 }n\text{位数字不出现 }0\},
\]
现在，首先把$[0,1]$作$10$等分，并舍去其中的$\left[0,\frac{1}{10}\right)$，即舍去第$1$个子区间，则在剩余区间中的点其十进位小数式中第一位小数就不会出现$0$了。其次，把剩余的$9$个子区间中的每一个再作$10$等分，并再舍去第$1$个子区间，$\cdots\cdots$如此进行下去，就得出$E_n$。易知，$m(E)$就是舍去的子区间之总长度为
\begin{align*}
m\left( E \right) =m\left( \left[ 0,1 \right] \backslash \bigcup_{n=1}^{\infty}{E_n} \right) =\frac{1}{10}+\frac{9}{10^2}+\frac{9^2}{10^3}+\cdots =\frac{1}{10}\sum_{n=0}^{\infty}{\frac{9^n}{10^n}}=1,
\end{align*}

则
\begin{align*}
m(E) &= \frac{1}{10} + \frac{9}{10^2} + \frac{9^2}{10^3} + \cdots = 1.
\end{align*}

\item 作点集（记$k=0,1,2,\cdots,9$）
\begin{align*}
E_k &= \{x\in[0,1]: \text{在十进位小数表示式 }x=0.a_1a_2\cdots\text{中的所有 }a_i\text{都不出现数字 }k\},
\end{align*}
易知，$E=\bigcup\limits_{k=0}^{9}E_k$，$E_k\sim[0,1]$（$k=0,1,2,\cdots,9$）.

现在，首先把$[0,1]$作$10$等分，并舍去其中的$\left[\frac{k}{10},\frac{k+1}{10}\right)$（$k=0,1,2,\cdots,9$），即舍去第$k+1$个子区间，则在剩余区间中的点其十进位小数式中第一位小数就不会出现$k$了。其次，把剩余的$9$个子区间中的每一个再作$10$等分，并再舍去第$k+1$个子区间，$\cdots\cdots$如此进行下去，就得出$E_k$。易知，舍去的子区间之总长度为
\begin{align*}
\frac{1}{10}+\frac{9}{10^2}+\frac{9^2}{10^3}+\cdots =\frac{1}{10}\sum_{n=0}^{\infty}{\frac{9^n}{10^n}}=1,
\end{align*}
由此可得$m(E_k)=0$.

\item 令$E_j=\{x\in(0,1): x\text{的小数展开式中不含数字 }j\}$（$j=0,1,\cdots,9$），则$E=\bigcup\limits_{j=0}^{9}E_j$。注意到把$[0,1]$中的数以九进位小数展开时，就可知$E_j$与$[0,1]$是对等的，即$\overline{E_j}=c$（$j=0,1,2,\cdots,9$）.

此外，将$[0,1]$作$10$等分，并舍去区间$[\frac{k}{10},\frac{k+1}{10})$（$k=0,1,2,\cdots,9$）的第$j+1$个子区间，则在余下的点的小数表示式中不出现（第一位）$j$。再将余下的区间又各$10$等分，又舍去第$j+1$个子区间,$\cdots\cdots$如此进行下去，就留下了$E_j$。易知共舍去之区间总长为
\begin{align*}
\frac{1}{10}+\frac{9}{10^2}+\frac{9^2}{10^3}+\cdots =\frac{1}{10}\sum_{n=0}^{\infty}{\frac{9^n}{10^n}}=1,
\end{align*}
即$m(E_j)=0$，$m(E)=0$.

\item \begin{enumerate}[(i)]
\item 设$\{x_k\}\subseteq E$：$x_k\to x$（$k\to\infty$），且令
\begin{align*}
x = \sum\limits_{n=1}^{\infty}\frac{b_n}{10^n}, \quad \text{若 }|x_k - x| < \frac{1}{10^p}, \text{则 }b_p = 2\text{ 或 }7.
\end{align*}
故$x\in E$，即$E$是闭集.

\item $\overline{E}=c$，$E$不是可数集.

\item 由$E$是闭集以及$E\neq[0,1]$，可知$E$不在$[0,1]$中稠密.

\item $E$是闭集,进而是可测集.将$[0,1]$作$10$等分，并舍去除第$3$和$8$个子区间即$[\frac{2}{10},\frac{3}{10})$和$[\frac{7}{10},\frac{8}{10})$之外的区间，则在余下的点的小数表示式中第一位只会是$2$或$7$.再将余下的区间又各$10$等分,又舍去除第$3$和$8$个子区间之外的区间,$\cdots \cdots$,如此进行下去,就留下了$E$。易知舍去之区间总长为
\begin{align*}
\frac{8}{10}+\frac{8}{10}\cdot \frac{2}{10}+\frac{8}{10}\cdot \frac{2^2}{10^2}+\cdots =\frac{8}{10}\sum_{n=0}^{\infty}{\left( \frac{2}{10} \right) ^m}=1.
\end{align*}
故
\begin{align*}
m(E) = 1 - \frac{8}{10}\sum_{n=0}^{\infty}{\left( \frac{2}{10} \right) ^m}= 0.
\end{align*}
\end{enumerate}
\end{enumerate}

\end{solution}

\begin{example}
试证明下列命题：
\begin{enumerate}[(1)]
\item 设$E\subseteq\mathbb{R}$，且存在$q$：$0<q<1$，使得对任一区间$(a,b)$，都有开区间列$\{I_n\}$：
\begin{align*}
E\cap(a,b)\subseteq\bigcup\limits_{n=1}^{\infty}I_n,\quad \sum\limits_{n=1}^{\infty}m(I_k)<(b-a)q.
\end{align*}
则$m(E)=0$.

\item 设$\alpha>2$，作$\mathbb{R}$中点集：
\[
E = \{x: \text{存在无限个分数 }\frac{p}{q},p\text{ 与 }q\text{ 是互素的自然数}, \text{使得 }|x-\frac{p}{q}|<\frac{1}{q^\alpha}\},
\]
则$m(E)=0$.

\item 设$a_n>0$（$n\in\mathbb{N}$），$\sum\limits_{n=1}^{\infty}a_n<+\infty$，作点集
\[
E = \{x\in(0,1): \text{存在无限个既约分数 }\frac{p}{q},p\text{ 与 }q\text{ 是互素的自然数}, \text{使得 }|x-\frac{p}{q}|<\frac{a_q}{q}\},
\]
则$m(E)=0$。

\item 设$E\subseteq[0,1)$是可测集，且$m(E)>0$，令
\[
A = \{x\in[0,1): \text{存在 }n\in\mathbb{N}\text{ 以及 }t\in E,\text{使得 }x=\{nt\}\},
\]
（$\{nt\}$是$nt$的小数部分）则$m(A)=1$.
\end{enumerate}
\end{example}
\begin{proof}
\begin{enumerate}[(1)]
\item 因为$m(E)=m\left(E\cap\bigcup\limits_{n=1}^{\infty}(-n,n)\right)\leqslant\sum\limits_{n=1}^{\infty}m(E\cap(-n,n))$，所以只需指出对任意的$(a,b)\subseteq\mathbb{R}$，有$m^*(E\cap(a,b))=0$。

对任意的$(a,b)\subseteq\mathbb{R}$，由题设知，存在$(a_n,b_n)$（$n\in\mathbb{N}$），使得
\begin{align}
\bigcup\limits_{n=1}^{\infty}(a_n,b_n)\supseteq E\cap(a,b),\text{且}\sum\limits_{n=1}^{\infty}(b_n-a_n)<q(b-a).\label{eq:::-290hrewfgesr64uj67urhds3ew21.1}
\end{align}
从而
\begin{align*}
m(E\cap(a,b))\leqslant m\left(\bigcup\limits_{n=1}^{\infty}(a_n,b_n)\right)\leqslant\sum\limits_{n=1}^{\infty}(b_n-a_n)<q(b-a).
\end{align*}
对每个区间$(a_n,b_n)$（$n\in\mathbb{N}$），由题设知，存在$(a_{nk},b_{nk})$（$k\in\mathbb{N}$），使得
\begin{align}
\bigcup\limits_{k=1}^{\infty}(a_{nk},b_{nk})\supseteq E\cap(a_n,b_n),\text{且}\sum\limits_{k=1}^{\infty}(b_{nk}-a_{nk})<q(b_n-a_n),\quad \forall n\in\mathbb{N}.\label{eq:::-290hrewfgesr64uj67urhds3ew21.2}
\end{align}
从而由\eqref{eq:::-290hrewfgesr64uj67urhds3ew21.1}\eqref{eq:::-290hrewfgesr64uj67urhds3ew21.2}式可知
\begin{align*}
m(E\cap(a,b))&\leqslant m\left(\bigcup\limits_{n=1}^{\infty}(a_n,b_n)\right)\leqslant m\left(\bigcup\limits_{n=1}^{\infty}\bigcup\limits_{k=1}^{\infty}(a_{nk},b_{nk})\right)\\
&\leqslant\sum\limits_{n=1}^{\infty}\sum\limits_{k=1}^{\infty}(b_{nk}-a_{nk})<\sum\limits_{n=1}^{\infty}q(b_n-a_n)<q^2(b-a).
\end{align*}
依此类推，继续做下去可得，对$\forall k\in\mathbb{N}$，都有
\begin{align*}
m(E\cap(a,b))<q(b-a)<q^2(b-a)<\cdots<q^k(b-a).
\end{align*}
令$k\rightarrow+\infty$，得$m(E\cap(a,b))=0$。再由$(a,b)$的任意性可得
\begin{align*}
m(E)=m\left(E\cap\bigcup\limits_{n=1}^{\infty}(-n,n)\right)\leqslant\sum\limits_{n=1}^{\infty}m(E\cap(-n,n))=0.
\end{align*}

\item 只需证$m(E_0)=m(E\cap(0,1])=0$。因为由\hyperref[theorem:外测度的平移不变性]{测度的平移不变性}知
\begin{align*}
m\left( E \right) =m\left( E\cap \bigcup_{n=-\infty}^{+\infty}{\left( n,n+1 \right]} \right) \leqslant \sum_{n=-\infty}^{+\infty}{m\left( E\cap \left( n,n+1 \right] \right)}=\sum_{n=-\infty}^{+\infty}{m\left( \left\{ n \right\} +E\cap \left( 0,1 \right] \right)}=\sum_{n=-\infty}^{+\infty}{m\left( E\cap \left( 0,1 \right] \right)}.
\end{align*}
记在$p,q$固定时$E_0$中之点集为$E_{p,q}$，又记$E_q=\bigcup\limits_{p=1}^{q}E_{p,q}$，由$E_0$的定义知,当$x\in E_0$时,存在无穷多个$p,q$，使得$x\in E_{p,q}$.而任意的$E_q$中只包含有限个$E_{p,q}$,故也存在无穷多个$q$,使得$x\in E_q$.故由\refthe{Set Theory-theorem:上、下极限集的刻画}知$x\in\varlimsup\limits_{q\to\infty}E_q$，由此知$E\subseteq\varlimsup\limits_{q\to\infty}E_q$.

此外，由$m(E_{p,q})=\frac{2}{q^{\alpha}}$可知，$m(E_q)\leqslant \frac{2q}{q^\alpha}=\frac{2}{q^{\alpha-1}}$。从而我们有$\sum\limits_{q=1}^{\infty}m(E_q)<+\infty$。因此由\rrefthe{theorem:Borel-Cantelli引理}{theorem:Borel-Cantelli引理-1}知$m\left(\varlimsup\limits_{q\to\infty}E_q\right)=0$，也就有$m(E)=0$.

\item 考查$E_q=\bigcup\limits_{p=0}^{q}(\frac{p-a_q}{q},\frac{p+a_q}{q})$（$q\in\mathbb{N}$，易知$E\subseteq\bigcup\limits_{q\geqslant m}E_q$（任意$m\in\mathbb{N}$）。
\end{enumerate}
\begin{enumerate}[(1)]
\setcounter{enumi}{3}
\item 由(外)测度的定义知,对任给$\varepsilon>0$，存在开区间列$\{I_k\}$，使得$\bigcup\limits_{k=1}^{\infty}I_k\supseteq E$，$\sum\limits_{k=1}^{\infty}m(I_k)<\frac{m\left( E \right)}{1-\frac{\varepsilon}{2}}$。从而存在$I_{k_0}$,
\begin{align*}
m(I_{k_0}\cap E)>(1-\frac{\varepsilon}{2})m(I_{k_0}).
\end{align*}
取有理数$r_0,r_1$，使$J\triangleq(r_0,r_1)\supset I_{k_0}$，且$m(J)<\frac{m\left( I_{k_0} \right)}{1-\frac{\varepsilon}{2}}$，则有
\begin{align*}
m(J\cap E)\geqslant m(I_{k_0}\cap E)>(1-\varepsilon)\cdot m(J).
\end{align*}
令$x_0=\frac{a}{n},x_1=\frac{b}{n},J_l=(\frac{l}{n},\frac{1+l}{n})$，其中$l$满足：$a\leqslant l\leqslant b-1$，使得$m(J_l\cap E)\geqslant(1-\varepsilon)m(J_l)$。因此，记$H=\{x\in[0,1):\frac{x}{n}\in J_l\cap E\}$，则$m(H)\geqslant1-\varepsilon$。注意到$H\subseteq A$，故$m(A)\geqslant1-\varepsilon$。由$\varepsilon$的任意性，知$m(A)=1$.
\end{enumerate}
\end{proof}


















\end{document}